% !PRINTED 97
\weberpage{97}
%<<CONT
und dafür können wir mit Benutzung der Bezeichnung (6) setzen

\weq{8}{(\varrho, \eta) = \overset{n}{\Sigma}\, \varrho^{\sigma} r^n,}

worin $n$ alle Zahlen der Gruppe $\mathfrak{N}$ durchläuft, und in $\sigma$ für
$\nu_1, \nu_2 \ldots \nu_\mu$ jedesmal die Indices von $n$ zu setzen sind.

Es mögen nun mit $r_1, r_2, \ldots r_\mu$ primitive Einheitswurzeln der Grade
$q_1, q_2, \ldots q_\mu$ bezeichnet sein, und

\[ r = r_1 r_2, \ldots r_\mu \]

gesetzt werden. Dann lässt sich der Ausdruck (8) in ein Product zerlegen,
nämlich, wenn wir

\[ n \equiv n_h \ (\mathrm{mod}\ q_h) \]

setzen, in:

\weq{9}{(\varrho, \eta) = \Sigma\, \varrho^{\varepsilon_1 \nu_1} r_1^{n_1}\, \Sigma\, \varrho^{\varepsilon_2 \nu_2} r_2^{n_2} \ldots \Sigma\, r_\mu^{\varepsilon_\mu \nu_\mu} r_\mu^{n_\mu}.} % [sic: Weber prints r_\mu as the final first base; the parallel factors require \varrho.]

Die Factoren dieses Productes sind nun genau die Resolventen, die wir in
\S.\,172 des ersten Bandes untersucht haben, und wir erhalten, wenn wir

\weq{10}{\overset{n_1}{\Sigma}\, \varrho^{\nu_h} r^{n_h} = (\varrho, \eta_h), \quad h = 1, 2, \ldots \mu} % [sic: Weber prints the upper bound n_1 and the base r^{n_h}; the general h factor requires n_h and r_h^{n_h}.]

setzen:

\weq{11}{(\varrho, \eta) = (\varrho^{\varepsilon_1}, \eta_1)(\varrho^{\varepsilon_2}, \eta_2) \ldots (\varrho^{\varepsilon_\mu}, \eta_\mu).}

Wenn wir hierauf die Formeln (5), (6), (25), (26) des eben angeführten
Paragraphen anwenden, so erhalten wir:

\weq{12}{(\varrho, \eta)(\varrho^{-1}, \eta) = m,}

\weq{13}{(\varrho, \eta)^3 = m(a + b\varrho), \quad (\varrho^2, \eta)^3 = m(a + b\varrho^2),}

worin $a$ und $b$ gewisse ganze Zahlen sind, die sich aus den genannten
Formeln berechnen lassen, die aber nach der Wahl der $\varepsilon$
verschieden ausfallen.

Aus (12) und (13) folgt dann noch:

\weq{14}{m = a^2 - ab + b^2.}

Hiernach können wir die cubische Gleichung aufstellen, deren Wurzeln die
Grössen $\eta, \eta', \eta''$ sind. Sie möge lauten:

\weq{15}{\eta^3 - \alpha \eta^2 + \beta \eta - \gamma = 0,}

worin $\alpha, \beta, \gamma$ ganze Zahlen sind. Was zunächst $\alpha$
betrifft, so ist es gleich der Summe aller primitiven $m^{\mathrm{ten}}$
Einheitswurzeln, also

\[ \alpha = \Sigma\, r_1\, \Sigma\, r_2 \ldots \Sigma\, r_\mu, \]

und dieser Werth ist $= 0$ oder $= (-1)^{\mu}$,
je nachdem 9 unter den Factoren von $m$ vorkommt oder nicht. Die beiden
anderen Co\"efficienten lassen sich ganz so berechnen, wie an der er-
%CONT>>
